% =============================================================================
% THE UFFIZI GUIDE
% A plain-language tour of every algorithm in the project: how each one works,
% how they differ and why they produce the numbers we actually observe.
% Build: pdflatex uffizi_guide.tex  (twice, for the table of contents)
% =============================================================================
\documentclass[11pt]{article}

\usepackage[a4paper,margin=2.4cm]{geometry}
\usepackage[T1]{fontenc}
\usepackage{lmodern}
\usepackage{amsmath,amssymb,mathtools}
\usepackage{xcolor}
\usepackage{booktabs}
\usepackage{array}
\usepackage{enumitem}
\usepackage[hidelinks]{hyperref}

\definecolor{inkblue}{RGB}{31,58,95}
\definecolor{accent}{RGB}{192,85,43}
\newcommand{\key}[1]{\textit{#1}}            % first use of a term
\newcommand{\why}[1]{\textcolor{inkblue}{#1}}
\setlength{\parskip}{0.45em}
\setlength{\parindent}{0pt}
\setlength{\emergencystretch}{2em}  % absorb tiny justification overflows

\title{\textbf{The Uffizi Guide}\\[0.3em]
\large A plain-language tour of the algorithms:\\
how they work, how they differ and why they produce the numbers they do}
\author{Uffizi RL project}
\date{}

\begin{document}
\maketitle
\thispagestyle{empty}

\begin{abstract}
\noindent This guide assumes you know nothing about reinforcement learning. It starts
from the everyday idea of learning by trial and error, builds up the vocabulary one
piece at a time and only then writes down an equation, explaining what every symbol is
doing. It covers all the methods the project uses: tabular Q-learning and Double
Q-learning on the toy museum; the handcrafted rules we compare against; MaskablePPO,
plain PPO and DQN on the full museum; and the training and evaluation scaffolding
(action masking, the curriculum, common random numbers, multiple seeds). The last third
of the guide does something most write-ups skip: it explains, method by method, \key{why}
the numbers came out the way they did, so that you can look at our tables and figures and
say not just \emph{what} happened but \emph{why} it had to happen.
\end{abstract}

\tableofcontents
\newpage

% =============================================================================
\section{The problem, in one breath}
% =============================================================================

The Uffizi Gallery in Florence takes roughly five thousand visitors a day through about
a hundred rooms. Almost all of them want the same four things: the two Botticelli
paintings, the Leonardo and the Raphael. So those four rooms are jammed shoulder to
shoulder while the rest of the museum sits nearly empty. That is the whole problem in a
sentence: not too many people for the building, but too many people for four rooms at
the same moments.

The project attacks this on two fronts that, for most of the document, you should keep
firmly apart.

The first front is the \key{museum-wide simulation}. We build an artificial crowd of
five thousand people, send them through the museum minute by minute and ask what
happens to congestion and to visitor satisfaction if the management changes the rules
(timed entry, booked slots for the masterpieces, longer hours and so on). That front
does not use any of the learning algorithms in this guide; it is a simulation of many
simple people, not one clever one.

The second front and the one this guide is about, is the \key{single optimal visitor}.
Take the redesigned museum as fixed and ask: what should one smart visitor actually
\emph{do}? When should they book? Which masterpieces? How should they pace the day so
they arrive at each painting at the right time and still get out before closing? That is
a sequence of decisions made under uncertainty, with consequences that only show up
later. That kind of problem is exactly what reinforcement learning is for and it is why
the algorithms below exist in the project at all.

\paragraph{How the two fronts meet and why we keep them apart.} The two fronts connect at a
single gate. The simulation decides whether a set of management changes is even acceptable: a
director will adopt a redesign only if it raises visitor welfare \emph{without} collapsing
ticket revenue, so the simulation's job is to find a bundle of changes (booked masterpiece
slots, longer hours, gentle pricing and so on) that passes that test. Only once a redesign
passes does the second front open: take that redesigned museum as fixed and ask what a
single optimal visitor does inside it. So the learning algorithms in this guide never touch
the question ``is the redesign worth doing?''; they answer the sharper question ``given that
we did it, how should a smart visitor use it?'' Keeping the two apart is what stops us from
confusing an \emph{average} over a whole messy population (what the simulation reports) with
the \emph{best} behaviour of one ideal visitor (what the agent learns). Those are different
questions with different answers and much of the confusion in this kind of project comes
from running them together.

% =============================================================================
\section{The vocabulary of reinforcement learning}\label{sec:vocab}
% =============================================================================

Before any algorithm, we need six words. We will introduce them with the museum visitor
as the running example, because every one of these words has a concrete meaning here.

\paragraph{State.} The \key{state} is everything the decision-maker knows at the moment
it has to act. For our visitor, a state might be: ``I am in room A11, it is 11:40, this
room is crowded and I have already seen Botticelli and Leonardo.'' The state is a
snapshot. Crucially, a good state contains \emph{everything that matters for the next
decision} and nothing the decision cannot depend on. We will see later that getting the
state wrong, leaving out something the right action secretly depends on, quietly breaks
the learning.

\paragraph{Action.} The \key{action} is the choice available in a state: ``move to the
next room,'' ``stay another minute,'' or in the booking problem, ``reserve the Leonardo
slot for 35 days from now.'' The set of legal actions can depend on the state (you can
only walk to a room that is physically next door) and that simple fact will turn out to
matter enormously.

\paragraph{Reward.} The \key{reward} is a single number the world hands back after each
action, measuring how good that step was right now. In the museum it is built from the
art you appreciated this minute, discounted if the room was crowded, minus a small cost
for time spent. Reward is local and immediate. It is \emph{not} the same as success;
success is about the whole visit and the reward only scores one step.

\paragraph{Return.} The \key{return} is the total reward summed over the entire visit,
from now until you walk out the door. This is what we actually care about. The whole
difficulty of the field is that you choose actions one at a time, getting only immediate
rewards, but you are graded on the return, which depends on every future step too.

\paragraph{Discount.} Because the future is uncertain and a payoff now is worth a little
more than the same payoff later, we add up the rewards with a \key{discount factor}
$\gamma$ between $0$ and $1$. A reward $k$ steps in the future is multiplied by
$\gamma^{k}$. The discounted return from time $t$ is
\begin{equation}
G_t \;=\; r_t \;+\; \gamma\, r_{t+1} \;+\; \gamma^2 r_{t+2} \;+\;\cdots\;=\; \sum_{k=0}^{\infty}\gamma^{k}\, r_{t+k}.
\label{eq:return}
\end{equation}
Read equation~\eqref{eq:return} slowly. $G_t$ is the return starting now. $r_t$ is the
reward we get this step, counted in full. The next reward $r_{t+1}$ is multiplied by
$\gamma$, the one after by $\gamma^2$ and so on, so rewards far in the future are shrunk
toward zero. If $\gamma$ is close to $1$ (we use $0.99$ and higher), the agent is
\key{far-sighted}: a reward forty steps away still counts for a lot, which is essential
here because the payoff for a good booking can land hundreds of minutes later. If
$\gamma$ were small, the agent would be a short-term creature that ignores the future.

\paragraph{A worked example of the return.} Numbers make this concrete. Suppose a short
visit hands out rewards $10, 0, 6, 4$ over four minutes and then ends. The plain
(undiscounted) total is $20$. With $\gamma = 0.99$ the return is
$10 + 0.99\cdot 0 + 0.99^2\cdot 6 + 0.99^3\cdot 4 = 10 + 0 + 5.88 + 3.88 = 19.76$, only a
hair under $20$ because a discount near one barely shrinks anything. Now redo the sum with
a short-sighted $\gamma = 0.5$: $10 + 0 + 0.25\cdot 6 + 0.125\cdot 4 = 12.0$. The same
visit is ``worth'' $19.76$ to a patient agent and only $12.0$ to an impatient one,
because the impatient one almost throws away the later $6$ and $4$. This is the entire
reason we set $\gamma$ as high as $0.999$ for the booking agent: the reward for booking
well does not arrive until check-in, possibly four hundred minutes later and at
$\gamma = 0.999$ a reward four hundred steps out is still worth $0.999^{400}\approx 0.67$
of its face value, whereas at $\gamma = 0.99$ it would be worth $0.99^{400}\approx 0.018$,
effectively invisible. The discount is not a technicality; it decides whether the agent
can even feel the consequence of its most important choice.

\paragraph{Policy and value.} A \key{policy}, written $\pi$, is the agent's strategy: a
rule that, given a state, says which action to take (or gives a probability to each
action). Learning is nothing more than improving the policy. To improve it we need a way
to score situations and that is the \key{value}. The \key{action-value} $Q(s,a)$ is the
return we expect if we take action $a$ in state $s$ and then act well forever after:
\begin{equation}
Q^{\pi}(s,a) \;=\; \mathbb{E}\!\left[\,G_t \;\middle|\; s_t = s,\; a_t = a,\; \text{then follow } \pi \right].
\label{eq:qdef}
\end{equation}
The symbol $\mathbb{E}[\cdot]$ means ``the average over all the random things that could
happen.'' So $Q(s,a)$ answers a very practical question: ``if I do $a$ here and play well
afterward, how much total reward should I expect?'' If we knew $Q$ exactly, the best
policy would be trivial: in every state, take the action with the largest $Q$. Every
algorithm in this guide is, at heart, a different way of getting at $Q$ or at the policy
directly.

\paragraph{The Markov decision process.} Put these together and you have a \key{Markov
decision process} (MDP): states, actions, a reward for each step, a rule for how the
state changes and a discount. ``Markov'' just means the state is a complete enough
summary that the future depends only on the present state and action, not on the whole
history of how you got there. When we later say a model ``broke because the task was not
Markov in the state it could see,'' we mean exactly this property was violated: the right
action depended on something missing from the state.

\paragraph{One step of the visitor's life, start to finish.} It helps to watch a single
turn of the loop, because every algorithm in this guide is just a different way of
improving one piece of it. The visitor is in a state $s$ (room A11, 11:40, crowded,
Botticelli already seen). The policy looks at $s$ and chooses an action $a$ (say, move on
to the corridor). The world responds: it hands back a reward $r$ for that step (a little
appreciation banked, discounted because the room was crowded, minus the time cost) and
moves the visitor to a new state $s'$ (in the corridor, 11:41). Now the loop repeats from
$s'$. The visit is a long chain of these $(s, a, r, s')$ transitions and the agent's only
job is to make choices so that the \emph{sum} of all those little rewards, the return, is
as large as possible. Value-based methods (Q-learning, DQN) attack this by learning the
number $Q(s,a)$ for each move and then always picking the largest; policy-gradient methods
(PPO) attack it by directly nudging the probabilities of the moves up or down. Keep this
one-step picture in your head; everything below is a variation on it.

% =============================================================================
\section{The five elements, made concrete for each algorithm}\label{sec:concrete}
% =============================================================================

A natural question, having met the six words, is: ``so which exactly is the state in
Q-learning? And in PPO?'' The clean answer carries an important idea:
\emph{state, action, reward, return and discount belong to the problem, the MDP, not to the
algorithm.} An algorithm is a recipe for learning a good policy \emph{within} a given MDP; it
does not change what the state or the reward is. So the right way to read the question is by
\emph{setting} and there are three. The first is the \key{toy museum} (tabular Q-learning and
Double Q-learning). The other two both live on the \emph{full} museum and are easy to confuse,
so we keep them apart. The \key{as-is museum} (no interventions) is the museum as it runs
today: the visitor simply \emph{walks} a planned route through the crowds and \emph{books
nothing}. The \key{RAMA museum} (the intervention) gates the masterpieces behind booked
time-slots, so there \emph{the decision is the booking}. The deep methods are compared across
both: \emph{PPO and DQN appear twice}, once as a no-booking walk on the as-is museum (the
baselines) and once as a booking agent on the RAMA museum, while \emph{MaskablePPO appears
only on the RAMA museum}. Within each setting the five elements are identical across the
algorithms that share it; what changes is only how the state is \emph{represented} and how the
action is \emph{chosen}.

\subsection{The toy MDP (Q-learning and Double Q-learning)}

These two algorithms face the \emph{same} MDP; nothing in this list differs between them.
\begin{description}[leftmargin=0pt,style=nextline,itemsep=2pt]
\item[State] A discrete five-part label: (current room; time-slice of the day, one of ten;
Botticelli crowd level, one of four; Leonardo crowd level, one of four; a bitmask of which
masterpieces have been seen). About $123{,}000$ distinct states, stored in a table.
\item[Action] From the current room: \emph{stay}, or \emph{move} to one of the physically
adjacent rooms. Only legal moves are offered.
\item[Reward] Per step,
$\text{importance}\times\text{novelty}/(1+\alpha_A\,\text{density}^2)-\text{step cost}$,
with $\alpha_A = 6$ and step cost $0.5$, plus a closing-time pressure, an exit bonus ($+15$),
and a no-exit penalty ($-50$).
\item[Return] The discounted sum of those per-step rewards over the whole visit,
$G_t = \sum_k \gamma^k r_{t+k}$ (equation~\eqref{eq:return}).
\item[Discount] $\gamma = 0.99$.
\end{description}

\paragraph{What the five-part label looks like.} At the entrance, opening time, the state is
\[
s=\bigl(\underbrace{\texttt{ENTRY}}_{\text{room}},\ \underbrace{0}_{\text{time-slice}},\ \underbrace{0}_{\text{Botticelli crowd}},\ \underbrace{0}_{\text{Leonardo crowd}},\ \underbrace{0}_{\text{seen mask}}\bigr):
\]
the opening slice, both masterpiece rooms empty, nothing seen yet. The crowd levels are what make
timing learnable. By mid-morning (time-slice $3$) the table reads Botticelli at its peak,
$(\dots,3,3,1,\dots)$: crowd level $3$ of $3$ at Botticelli while Leonardo is still light ($1$), so a
Q-learner that has learned \emph{temporal arbitrage} waits. Two slices later $(\dots,5,1,3,\dots)$ the
wave has rolled on: Botticelli has emptied to $1$ (enter now) while Leonardo has filled to $3$. The
last entry is a six-bit mask over the must-see rooms (\texttt{A11}, \texttt{A12}, \texttt{A16},
\texttt{A35}, \texttt{A38}, \texttt{E4}); after seeing only Botticelli's \texttt{A11} it reads $1$
(binary $000001$), so the agent can tell ``already seen'' from ``still to see''.
\emph{What differs between the two algorithms is only the bookkeeping.} Q-learning stores one
table $Q(s,a)$ and updates it with the single-table rule (equation~\eqref{eq:qupdate}),
choosing the action with the largest entry in the row. Double-Q stores \emph{two} tables and
updates them with the decorrelated rule (equation~\eqref{eq:doubleq}), choosing by the average
of the two. Same state, same actions, same reward, same return, same discount.

\paragraph{Where do the toy reward numbers come from?} They are hand-tuned \emph{shaping}
constants, not derived quantities and what matters is their \emph{relative} size rather than
the absolute values (multiplying every reward by the same factor leaves the optimal policy
unchanged). Each one buys a specific behaviour:
\begin{itemize}[leftmargin=1.4em,itemsep=2pt]
\item $\alpha_A = 6$ sets how sharply a crowd discounts a painting: in a full room
($\text{density}=1$) the value falls to $1/(1+6)\approx 14\%$ of its empty-room value. We made
it \emph{large} on purpose, so that visiting off-peak clearly beats visiting at the peak; the
whole point of the toy is to make temporal arbitrage worth learning. The full museum
deliberately uses a much \emph{gentler} discount, $\alpha=0.5$; a harsh one there produced the
``fake art lover'' that skipped crowded masterpieces, one of our documented failures.
\item the step cost $0.5$ is a small per-minute charge so dithering is never free: small enough
that a good room is still worth entering, large enough that loitering in a low-value room is
net-negative. It is what makes the agent move on.
\item the exit bonus $+15$ makes leaving worth real reward (about two masterpiece visits), so
the agent does not squat, while staying small enough that it does not abandon the art to bolt
for the door.
\item the no-exit penalty $-50$ is the closing-time guillotine: large enough to dominate a
whole episode's positive return ($\sim$80--100), so being locked in roughly wipes out the day,
but deliberately \emph{finite}. An infinite penalty (or zeroing all reward) would be flat
across policies and give no gradient, which is precisely the unlearnable-reward failure we
describe later.
\end{itemize}
So these were set by trial and error to produce sensible behaviour; several of them exist
\emph{because} an earlier value failed and the booking reward's weights were shaped the same
way. They are a hand-tuned part of the model, not first-principles constants; the
\emph{rankings} they produce (off-peak beats peak, leaving beats squatting) are robust to the
exact numbers, which is what we rely on.

\subsection{The full museum: as-is (no booking) versus RAMA (booking)}

On the full museum there are \emph{two} MDPs, not one and confusing them is the easiest
mistake to make. The first is what the deep \emph{baselines} (\texttt{PPO-base},
\texttt{DQN-base}) solve and it involves \emph{no booking at all}:
\begin{description}[leftmargin=0pt,style=nextline,itemsep=2pt]
\item[State] An eight-number summary of a visitor walking a fixed planned route (room-value
drain, importance, density, on-target flag, fraction of the plan done, time, egress slack,
distance to the next room).
\item[Action] Just two (\texttt{Discrete(2)}): \emph{stay} or \emph{move on}. Nothing is booked.
\item[Reward] The per-minute museum reward (appreciation under the crowd discount, closing
pressure, completion, step cost); no booking terms.
\item[Return / Discount] The discounted sum, with $\gamma = 0.995$ (art lover) or $0.997$ (tourist).
\end{description}
This is the museum \emph{as it runs today}: you walk in and take the masterpieces as you find
them, crowds and all. It is the fair opponent for the intervention.

\paragraph{What the eight numbers look like.} At the door (room \texttt{A1}, opening time) the
observation reads
\[
s=\bigl(\underbrace{0.00}_{\text{drain}},\ \underbrace{0.22}_{\text{value}},\ \underbrace{0.30}_{\text{density}},\ \underbrace{0}_{\text{on target?}},\ \underbrace{0.00}_{\text{plan done}},\ \underbrace{0.00}_{\text{time}},\ \underbrace{0.99}_{\text{egress}},\ \underbrace{0.10}_{\text{dist.\ next}}\bigr).
\]
Six moves later, standing in the first must-see room \texttt{A9}, the same eight slots read
\[ (0.21,\,0.78,\,0.02,\,1,\,0.00,\,0.01,\,0.98,\,0.00). \]
Taking the slots one at a time, with the two readings side by side:
\begin{description}[leftmargin=0pt,style=nextline,itemsep=3pt]
\item[1. Room-value drain \textmd{(\texttt{A1} $0.00$, \texttt{A9} $0.21$)}] How much of what the
\emph{current} room has to give you have already absorbed: the appreciation taken here so far divided
by the room's ideal dwell time, capped at $1$. $0$ means you just walked in and have taken nothing;
$1$ means you are sated and more minutes here are wasted; it climbs the longer you linger. At
\texttt{A9} you are one-fifth of the way through ($0.21$). This is the agent's ``have I seen enough of
this room?'' gauge; it is exactly the hidden progress variable whose absence once turned the visit
into a POMDP, because without it the agent cannot tell a fresh room from a drained one.
\item[2. Room value \textmd{(\texttt{A1} $0.22$, \texttt{A9} $0.78$)}] The cultural importance of the
room underfoot on the graph's $1$ to $10$ scale, divided by $10$. \texttt{A1} is a minor room ($2.2$);
\texttt{A9} is a masterpiece room ($7.8$). It tells the agent how much reward dwelling here can yield
at all, so it spends its minutes where they pay.
\item[3. Density \textmd{(\texttt{A1} $0.30$, \texttt{A9} $0.02$)}] The crowd in the current room right
now, measured as occupancy over capacity and clipped at $1$ ($0$ empty, $1$ at or over capacity). Here
\texttt{A9} happens to be nearly empty ($0.02$). Crowd discounts appreciation (the reward divides by
$1+\alpha\,\text{density}^2$), so this is the signal the agent uses to hit a masterpiece in its lull,
the temporal-arbitrage move.
\item[4. On-target flag \textmd{(\texttt{A1} $0$, \texttt{A9} $1$)}] $1$ when you are standing in the
\emph{target}, the next must-see your plan points to (the nearest must-see you have not yet finished),
$0$ otherwise. At \texttt{A1} you are not there ($0$); at \texttt{A9} you have arrived ($1$). It lets
the agent know when ``stay and appreciate'' is the right move rather than ``keep walking''.
\item[5. Fraction of the plan done \textmd{(\texttt{A1} $0.00$, \texttt{A9} $0.00$)}] How many of your
must-see rooms are fully appreciated, as a share of the whole plan: $0$ at the start, $1$ when every
must-see is done and only the exit remains. Still $0$ at \texttt{A9} because you have only just reached
the first one. It tells the agent how much of the visit's purpose is already banked.
\item[6. Time of day \textmd{(\texttt{A1} $0.00$, \texttt{A9} $0.01$)}] Minutes elapsed divided by the
length of the day: $0$ at opening, $1$ at closing. Crowds follow the clock and the museum shuts, so
every decision is time-sensitive.
\item[7. Egress slack \textmd{(\texttt{A1} $0.99$, \texttt{A9} $0.98$)}] Spare time to get out before
closing: (minutes left in the day minus the walking distance to the nearest exit) over the day length,
clipped to $[-1,1]$. Near $1$ is oceans of time; near $0$ means leave now; negative means you can no
longer make it out. This is the \emph{dense} ``head for the door'' signal we added after the
all-or-nothing exit penalty proved unlearnable.
\item[8. Distance to the next room \textmd{(\texttt{A1} $0.10$, \texttt{A9} $0.00$)}] How far the next
planned room is, in graph hops, divided by $20$ and capped at $1$. At \texttt{A1} the next must-see is
two hops off ($2/20=0.10$); at \texttt{A9} you are already in it, so $0$. It tells the agent how much
walking, time and crowd exposure stand between it and its next goal.
\end{description}

The second MDP is the \key{RAMA} task, what \texttt{MaskablePPO} solves (with unmasked PPO and
DQN as controls) and here \emph{the booking is the whole decision}:
\begin{description}[leftmargin=0pt,style=nextline,itemsep=2pt]
\item[State] A vector of about $25$ numbers: a booking-phase flag and step; the lead time
chosen so far; how many slots remain open at each of the four lead options; the time of day;
the three booked windows held; the three check-ins done; an ``access read'' of the current
room; the three frozen arrival estimates; and the current room's value, importance, density
and egress slack.
\item[Action] One of four (\texttt{Discrete(4)}): in the booking phase, the lead time
$\in\{1,7,21,35\}$ days and whether to take each masterpiece slot; in the visit phase,
\emph{stay} or \emph{move on}.
\item[Reward] The same per-minute museum reward plus the booking terms (the immediate off-pace
penalty and the no-show penalty).
\item[Return / Discount] The discounted sum over the two-phase episode, with $\gamma = 0.999$
(the booking payoff lands at check-in, hundreds of minutes later).
\end{description}

\paragraph{What the twenty-five numbers look like.} On the booking screen, before the visit and on
a busy day (\texttt{daily\_total}$=2500$), the art lover sees a vector that opens
\[
\bigl(\underbrace{1}_{\text{booking}},\ \underbrace{0}_{\text{step}},\ \underbrace{1.0}_{\text{lead so far}};\ \underbrace{0.12,\,0.38,\,0.62,\,1.0}_{\text{slots open at lead }1,7,21,35\text{ d}};\ \underbrace{0.0}_{\text{time}};\ \dots\bigr).
\]
Those four boxed numbers are the crux of the task: book only a day ahead and just \emph{12\%} of the
masterpiece windows are left; book a month ahead and \emph{all} of them are open. Here is the whole
vector, block by block, with the values from this opening state:
\begin{description}[leftmargin=0pt,style=nextline,itemsep=3pt]
\item[Where you are in booking \textmd{(slots 1 to 3: $1,\ 0,\ 1.0$)}] A \emph{phase flag} ($1$ = you
are on the booking screen, $0$ = you are walking the museum), a \emph{step} counter for which booking
sub-decision you are on (normalized, $0$ at the first), then the \emph{lead time chosen so far} as a
fraction of the $35$-day maximum ($1.0$ = the full month, before you have shortened it). Together they
tell the agent which decision it is making right now.
\item[The booking screen \textmd{(slots 4 to 7: $0.12,\ 0.38,\ 0.62,\ 1.0$)}] The fraction of
masterpiece slots still open if you reserve at each of the four lead times, $1$, $7$, $21$ and $35$
days ahead. On this busy day only $12\%$ remain a day before but all are free a month before, which is
the whole reason the agent learns to \emph{book earlier when the crowd is bigger}.
\item[Time of day \textmd{(slot 8: $0.0$)}] Minutes elapsed over the length of the day, $0$ here
because the visit has not started; it climbs to $1$ at closing once the agent is inside.
\item[Windows you hold \textmd{(slots 9 to 11: $0,\ 0,\ 0$)}] The start time (as a fraction of the
day) of the booked window for Botticelli, Leonardo and Raphael in turn, or $0$ if that masterpiece is
not booked. All zero here because nothing has been reserved on the opening screen.
\item[Check-ins done \textmd{(slots 12 to 14: $0,\ 0,\ 0$)}] One flag per masterpiece, $1$ once you
have actually entered its room inside its window and collected the reward, $0$ otherwise. All zero
before the visit.
\item[Access read of the room underfoot \textmd{(slots 15 to 18: $0,\ 0,\ 0,\ 0$)}] Four numbers that
describe the current room \emph{as a reservation}: is it a gated masterpiece, is it a masterpiece you
have not booked (so you cannot enter), may you appreciate it right now (your window is open or you are
checked in) and how many minutes until your window opens. All zero here because \texttt{A1} is an
ordinary room; they light up only when the agent stands in a masterpiece, connecting ``I get nothing
here'' to ``my window is not open yet''.
\item[Frozen pace prior \textmd{(slots 19 to 21: $0.23,\ 0.57,\ 0.66$)}] A running-average estimate,
learned on earlier episodes then frozen, of \emph{when in the day this visitor reaches} Botticelli,
Leonardo and Raphael. Here it expects Botticelli about a quarter into the day, Leonardo and Raphael
past the half. The agent books its windows around these times, so it reserves a slot it can actually
reach rather than one it will miss.
\item[The room-level tail \textmd{(slots 22 to 25: $0,\ 0.22,\ 0.05,\ 0.99$)}] The same four signals
as the as-is state: how drained the current room is, its value (importance over $10$), its crowd and
the egress slack to the exit. They let the single network handle the walking part of the episode once
the booking is settled.
\end{description}
\emph{What differs between the algorithms within the RAMA task} is only the machinery: DQN
outputs a value $Q(s,a)$ and acts greedily; PPO outputs action probabilities and samples;
MaskablePPO is PPO with illegal actions forced to zero before sampling
(equation~\eqref{eq:mask}).

\paragraph{A note on the real Uffizi.} Our ``booking'' is the \key{RAMA} intervention we
\emph{propose}: timed slots for the masterpiece rooms, which do not exist today. It should not
be confused with the entry-ticket booking that already exists, where a visitor either pre-books
admission (around 29 euros and skips the entrance queue) or buys at the door (around 25 euros
and queues). Our as-is baseline abstracts that entrance choice away and models only the
in-museum congestion; modelling the entry book-or-queue decision and its price gap would be a
natural extension.

\subsection{The whole thing, one line each}

\renewcommand{\arraystretch}{1.3}
\begin{center}\small
\setlength{\tabcolsep}{4.5pt}
\begin{tabular}{p{2.4cm} >{\raggedright\arraybackslash}p{3.4cm} >{\raggedright\arraybackslash}p{3.9cm} >{\raggedright\arraybackslash}p{2.7cm} l}
\toprule
\textbf{Algorithm} & \textbf{State} & \textbf{Action} & \textbf{Reward} & \textbf{$\gamma$} \\
\midrule
Q-learning & toy 5-tuple, one table & stay / move & toy per-step & 0.99 \\
Double-Q & toy 5-tuple, two tables & stay / move & toy per-step & 0.99 \\
PPO / DQN, as-is & 8-number walk & stay / move on & walk (no booking) & 0.995/0.997 \\
MaskablePPO, RAMA & 25-number booking & lead / book / pace, masked & booking per-step & 0.999 \\
PPO / DQN, RAMA & 25-number booking & lead / book / pace & booking per-step & 0.999 \\
\bottomrule
\end{tabular}
\end{center}

Reading across: the only things that ever change are which \emph{problem} the algorithm is
attached to (toy, the as-is walk, or RAMA booking) and, within a setting, how it stores the state and screens
the actions. That is the precise sense in which the five elements are properties of the setting; the algorithm is the method you bring to it.

% =============================================================================
\section{Learning from scratch: tabular Q-learning}
% =============================================================================

Start with the smallest version of the museum: a toy map of twelve rooms (entrance, two
corridors, the masterpieces, a couple of minor rooms, the exit). It is small enough that
we can keep a giant lookup table with one entry for every (state, action) pair and simply
fill in the numbers. This is \key{tabular} Q-learning and it is the cleanest way to see
the core idea before neural networks complicate the picture.

\paragraph{What the toy state actually is.} Concretely, the toy state is a five-part label:
which room you are in, which of ten time-slices of the day it is, how crowded Botticelli is
(quantised to four levels), how crowded Leonardo is (four levels) and a small bitmask of
which masterpieces you have already seen. Multiply those possibilities together and you get
on the order of a hundred and twenty thousand distinct states, each with a handful of moves.
That is small enough to keep every cell in memory and to visit each one many times over a
training run, which is the entire reason the toy exists: it lets us watch the pure algorithm
fill in a table and converge, with no neural network in the way to muddy what is happening.

\subsection{The idea before the equation}

Imagine a notebook with a row for every situation you could be in and a column for every
move and in each cell a running estimate of ``how good is this move, here, in the long
run?'' At the start the notebook is blank (all zeros), so the agent has no idea what is
good. It wanders, mostly at random and every time it takes a step it learns one tiny
lesson: ``the move I just made led to an immediate reward, plus a next situation that my
notebook says is worth such-and-such; so my old guess for that cell was a bit too high or
too low, let me nudge it.'' Repeat that nudge a few hundred thousand times and the
notebook fills in with remarkably accurate long-run values. That nudge is the entire
algorithm.

\subsection{The equation, piece by piece}

The one-step update, applied every time the agent is in state $s$, takes action $a$, gets
reward $r$ and lands in the next state $s'$, is
\begin{equation}
Q(s,a)\;\leftarrow\;Q(s,a)\;+\;\alpha\Big[\;\underbrace{r+\gamma\max_{a'}Q(s',a')}_{\text{the TD target: a better estimate}}\;-\;\underbrace{Q(s,a)}_{\text{old estimate}}\;\Big].
\label{eq:qupdate}
\end{equation}
Take it apart. $Q(s,a)$ on the left is the cell we are updating. The arrow means ``becomes.''
Inside the bracket is the \key{temporal-difference error}: the gap between a freshly
computed, better estimate and the stale one we had. The better estimate, called the
\key{TD target}, is $r + \gamma \max_{a'} Q(s',a')$: the reward we actually just received,
plus the discounted value of the best thing we could do from the next state $s'$. The
term $\max_{a'} Q(s',a')$ is the agent saying ``from where I landed, I will assume I play
optimally, so I take the largest value in that row.'' We subtract the old guess $Q(s,a)$
to get the error and $\alpha$ (the \key{learning rate}, we use $0.1$) controls how big a
nudge we make: small enough to be stable, large enough to learn in reasonable time.
This is called \key{bootstrapping}, because the estimate is updated toward a target that
itself contains an estimate ($Q(s',\cdot)$). We pull ourselves up by our own bootstraps,
and remarkably it converges.

\paragraph{Watching one cell learn.} Suppose the cell $Q(s,a)$ currently holds $0$ (we
know nothing yet), the agent takes $a$, receives reward $r = 5$ and lands in a state
$s'$ whose best known move is worth $\max_{a'}Q(s',a') = 20$. With $\gamma = 0.99$ and
$\alpha = 0.1$, the TD target is $5 + 0.99\cdot 20 = 24.8$, the TD error is
$24.8 - 0 = 24.8$ and the new value is $0 + 0.1\cdot 24.8 = 2.48$. We did not jump all the
way to $24.8$; we moved a tenth of the way, because $\alpha$ keeps each lesson small so
that one unlucky step cannot overwrite everything. The next time the agent visits the
same cell it nudges again and again and the value creeps toward its true long-run worth.
Notice what just happened conceptually: information about the good state $s'$ (worth $20$)
flowed \emph{backward} into the value of the earlier move $a$. Do this for every cell,
millions of times and value seeps backward from the exit and the masterpieces all the way
to the entrance, until the table encodes a coherent long-run plan even though every single
update only ever looked one step ahead.

\paragraph{Off-policy and why that word matters.} Q-learning is \key{off-policy}: the
$\max$ in the target assumes the agent will act \emph{optimally} from $s'$ onward, even
though, while exploring, it might actually take a random move next. In other words it
learns the value of the best policy while behaving according to a more adventurous one.
This is what lets it explore freely (try bad moves to learn they are bad) without
corrupting the thing it is learning. The contrast with PPO later, which is \key{on-policy}
and learns only about the policy it is currently running, is one of the deepest
distinctions in the field and it is the reason DQN can recycle old experience from a
memory buffer while PPO cannot.

\subsection{Two details that matter later}

\paragraph{Exploration.} If the agent always took what currently looks best, it would
lock onto the first decent path it found and never discover better ones. So it uses
\key{$\varepsilon$-greedy} action selection: with probability $\varepsilon$ it acts
randomly (explore), otherwise it takes the best-known move (exploit). We start at
$\varepsilon = 1.0$ (all exploration) and decay it multiplicatively toward $0.01$, so the
agent is adventurous early and disciplined late. This randomness is seeded and it is the
\emph{only} source of run-to-run difference on the deterministic toy, a fact we will need
when we explain why two seeds give slightly different numbers.

\paragraph{Masking the max.} A visitor can only step to an adjacent room. So in
equation~\eqref{eq:qupdate} the $\max_{a'}$ is taken only over \emph{legal} moves from
$s'$, never over impossible ones. This is a small, free piece of correctness in the
tabular world; it becomes a central design choice once we move to neural networks.

\subsection{What it learns and why}

On the toy map, the crowds are \key{deterministic waves}: Botticelli peaks at a fixed time
of day, Leonardo a bit later and so on, the same every run. The agent is never told this
schedule. Yet through equation~\eqref{eq:qupdate} it discovers \key{temporal arbitrage}:
visit a masterpiece in the valley before or after its crowd peak, not at the peak when a
naive plan would send you. \why{It learns this because the reward divides appreciation by
a crowd term: the same painting pays far more when the room is empty. Q-learning, summing
discounted rewards, sees that a small wait to enter off-peak buys a much larger reward,
and the bootstrapping propagates that future gain back to the earlier ``wait'' action.}
That is the whole point of the toy: to show that timing, learned from nothing but trial
and error, beats any fixed plan.

% =============================================================================
\section{The optimist's curse: Double Q-learning}\label{sec:doubleq}
% =============================================================================

Q-learning has a famous, subtle flaw and the project includes a second algorithm built
specifically to fix it. Understanding the flaw is worth a page because it also explains a
result that surprises people: on our toy, the ``improved'' algorithm does \emph{not} win.

\subsection{Why a single table is an optimist}

Look again at the TD target $r + \gamma \max_{a'} Q(s',a')$. While learning, the values in
the table are not yet correct; they are noisy estimates, some too high, some too low. Now
ask what the $\max$ does to noise. If several actions have the same true value but your
estimates are noisy, the $\max$ will, almost by definition, pick the one whose noise
happened to be \emph{positive}. So the maximum of noisy estimates is, on average, larger
than the maximum of the true values. Formally this is \key{Jensen's inequality}:
\begin{equation}
\mathbb{E}\Big[\max_{a'} Q(s',a')\Big] \;\ge\; \max_{a'} \mathbb{E}\big[Q(s',a')\big].
\label{eq:jensen}
\end{equation}
In words: the expected value of the maximum is at least the maximum of the expected
values. The agent's TD target is therefore systematically too optimistic. This is
\key{maximisation bias} or \key{overestimation}: the same noisy table is used both to
\emph{choose} the best next action and to \emph{score} it and that double duty inflates
the estimate.

\paragraph{A coin-flip example you can feel.} Imagine three moves whose true long-run
value is exactly $0$. Your estimates are noisy, so on a given step they might read
$+3, -2, +1$. The truth is that the best move is worth $0$, but $\max(+3, -2, +1) = +3$,
so your target says ``the best is worth $+3$.'' Another step the noise reads $-1, +4, 0$,
and the max is $+4$. Average the maxima over many steps and you get a number comfortably
above $0$, even though no move is actually worth more than $0$. You have convinced yourself
the future is rosier than it is, purely because you always quote the luckiest of several
noisy numbers. That is overestimation in one sentence and it compounds: an inflated value
for $s'$ flows backward through bootstrapping and inflates earlier cells too. Double-Q
breaks it because the table that \emph{spotted} the lucky $+4$ is not the table asked
\emph{how much it is really worth}; the second, independently noisy table will, on average,
quote something near the truth.

\subsection{The fix: two tables and a coin}\label{sec:doubleq-fix}

Double Q-learning (van Hasselt, 2010) breaks the self-flattery by keeping \emph{two}
tables, $Q_A$ and $Q_B$ and on each step flipping a coin to decide which one to update.
When updating $Q_A$, it uses $Q_A$ only to \emph{pick} the next action, but $Q_B$ to
\emph{score} it:
\begin{equation}
Q_A(s,a)\;\leftarrow\;Q_A(s,a)+\alpha\Big[\,r+\gamma\,Q_B\big(s',\,\textstyle\arg\max_{a'}Q_A(s',a')\big)-Q_A(s,a)\Big].
\label{eq:doubleq}
\end{equation}
The key is the inner part: $\arg\max_{a'}Q_A(s',a')$ asks $Q_A$ ``which action looks
best?'' and then $Q_B(s', \cdot)$ asks a \emph{different, independently noisy} table ``and
what is that action actually worth?'' Because the table that chooses is not the table that
scores, the lucky positive noise that fooled a single table no longer lines up and the
overestimation largely cancels. At decision time the agent acts on the average
$\tfrac12(Q_A + Q_B)$. The symmetric update (swap $A$ and $B$) runs when the coin lands the
other way.

\subsection{Why it does not help on our toy (and why that is the right result)}

Here is the result that looks wrong at first: across five random seeds, vanilla Q-learning
scores $83.3 \pm 3.8$ and Double-Q scores $83.1 \pm 3.7$. They are \key{statistically
tied}; the variant neither helps nor hurts. \why{The reason is that overestimation feeds
on noise and the toy has almost none. The crowd waves are deterministic and the rewards
are a fixed function of (state, action), so once the table converges there is no random
noise in the targets for the $\max$ to inflate. With no bias to remove, Double-Q's
machinery is solving a problem that is not there.} The small run-to-run wobble you see is
not target noise; it is exploration noise, which way the $\varepsilon$-greedy coin sent
the agent early in training and that affects both algorithms equally.

This is exactly why we report both. The real lesson is not ``Double-Q is better'' but
``Double-Q earns its keep precisely when the environment is noisy.'' Hold that thought:
when we get to DQN on the full, genuinely stochastic museum, the very same overestimation
will return and will hurt and the contrast will explain a real gap in the results.

% =============================================================================
\section{The rules we compare against: the handcrafted baselines}
% =============================================================================

To know whether learning is worth anything, we race the learned policy against five
\key{handcrafted} policies on the toy: fixed rules a sensible person might write by hand.
A learned policy that cannot beat ``follow the guidebook route'' has proven nothing.

\begin{itemize}[leftmargin=1.4em,itemsep=2pt]
\item \key{default\_path}: follow one fixed recommended route, the same for everyone.
\item \key{random}: at each step pick a legal move uniformly at random.
\item \key{greedy\_least\_crowded}: always step toward the emptiest neighbouring room.
\item \key{greedy\_value\_ratio}: always step toward the room with the best art-per-crowd ratio, right now.
\item \key{peak\_avoidance}: a hand-built timing heuristic that tries to dodge the known peaks.
\end{itemize}

These are not learning algorithms; they are fixed strategies, so they need no training.
They exist only as a yardstick. We will see in the results section that, strikingly, most
of them score \emph{negative}, worse than the random policy and the explanation of why a
``smart'' rule loses to coin-flipping is one of the more instructive parts of the story.

% =============================================================================
\section{When the table explodes: why we need neural networks}
% =============================================================================

The lookup table worked because the toy had only twelve rooms and a handful of state
features, about a hundred thousand cells in total, small enough to visit each many times
and fill in. The real museum destroys that. It has ninety-eight rooms, roughly a hundred
possible moves at each step, a continuous crowd density in every room and an observation
with several hundred numbers in it. A table over all of that would have more cells than we
could ever visit and almost every cell would stay blank forever. Worse, a table has no
notion that two similar situations should have similar values; it treats every cell as
unrelated, so nothing learned in one state helps in a nearby one.

\paragraph{What the full observation contains.} For the deep agent the state is a vector of
several hundred numbers, assembled fresh each minute: a one-hot code for the current room
(a $1$ in the slot for ``where I am,'' zeros elsewhere), the crowd density in every room and
whether that density is rising or falling, the time of day, an \key{egress slack} (how much
time is left minus how many hops to the nearest exit), a visited-flag for every room and,
importantly, a per-room \key{appreciation-progress} value recording how much of each room's
worth has already been banked. Several of these features earn their place only because
earlier versions of the agent broke without them, the egress slack and the
appreciation-progress especially, which is a story the failures section tells. The point for
now is that the network reads this whole vector and, because it is a smooth function, can
treat two similar vectors similarly, the generalisation a table can never manage.

The fix is \key{function approximation}. Instead of storing a separate number for every
situation, we train a \key{neural network} to compute values or action-probabilities from
the state. A network is just a flexible function with adjustable knobs (its weights): you
feed in the state numbers, it multiplies and adds and bends them through a few layers and
out comes either a value or a probability for each action. Because it is a smooth function,
it \key{generalises}: nudging the weights to be right about one crowded masterpiece room
automatically makes it more right about other crowded rooms. We use small networks, two
hidden layers of 128 or 256 units, which is plenty for this problem.

\paragraph{How a network learns, in plain terms.} A neural network is a stack of
\key{layers}. Each layer takes the numbers from the previous layer, multiplies them by its
weights, adds them up and passes the result through a simple bend (a nonlinearity) so the
whole thing can represent curves and not just straight lines. The first layer reads the
state; the last layer outputs what we want (a value, or a probability per action). Training
means adjusting the weights so the outputs are good. We measure ``how wrong'' the network
is with a \key{loss}, a single number that is large when the output is far from what it
should be (for DQN the loss is the squared TD error; for PPO it is the clipped objective,
flipped to a loss). Then we use \key{gradient descent}: compute, for every weight, the
direction that would most reduce the loss (the \key{gradient}) and take a small step that
way. Repeat on batch after batch of experience and the network slowly slides into weights
that produce good outputs. So ``the agent learns'' really means ``a loss is repeatedly
nudged downhill by tiny weight adjustments.'' Nothing more mystical than that is happening,
and keeping this picture in mind demystifies every training curve you will ever see: it is
just a loss going down (and the corresponding reward going up) as the steps accumulate. Everything from here
on replaces ``look up the cell'' with ``ask the network,'' but the underlying ideas, value,
return, temporal-difference error, are exactly the ones you already have.

% =============================================================================
\section{Learning the policy directly: PPO}
% =============================================================================

There are two broad ways to use a network. One is to keep estimating $Q$ and act greedily,
the neural descendant of Q-learning; that is DQN and we cover it shortly. The other is to
skip values as the main object and adjust the \emph{policy} itself, the probabilities of
each action, so that good actions become more likely. That second family is
\key{policy-gradient} and our workhorse, \key{PPO} (Proximal Policy Optimisation), lives
there. PPO is the method behind the project's headline results, so it is worth
understanding properly.

\subsection{Policy gradient, the core idea}

Forget values for a moment. Suppose the policy is a dial-box of probabilities: in each
state it assigns some chance to each action and we can turn the dials. The most direct
imaginable learning rule is: run a whole visit, see how much total reward it earned and
then turn the dials so that the actions taken on a \emph{good} visit become more likely and
the actions taken on a \emph{bad} visit become less likely. That is literally the original
policy-gradient method, called \key{REINFORCE} and its update has the shape
\begin{equation}
\theta \;\leftarrow\; \theta \;+\; \eta \, G \, \nabla_\theta \log \pi_\theta(a\mid s),
\label{eq:reinforce}
\end{equation}
where $G$ is the return the episode earned and $\nabla_\theta \log \pi_\theta(a\mid s)$ is
the direction in weight-space that makes action $a$ more probable. Multiplying by $G$ means:
if the episode was good ($G$ large and positive), step in the direction that increases the
probabilities of the actions you took; if it was bad, step the other way. Over many
episodes the dials drift toward whatever earns reward. No values, no $\max$, just ``do more
of what worked.''

The trouble with the raw recipe is \key{variance}. Using the whole-episode return $G$ to
judge \emph{every} action in that episode is terribly noisy: a single great booking can be
buried inside an otherwise mediocre visit and a clumsy move can ride the coat-tails of a
lucky day. The learning signal is unbiased on average but jumps around so much that training
is painfully slow. Two ideas tame it and together they turn REINFORCE into a method that
actually works at scale. First, instead of judging an action by the whole return, judge it
by how much better things went than \emph{expected}, which requires something that predicts
the expectation. Second, take small, controlled steps so the noisy signal cannot blow the
policy up on any single batch. The first idea gives us the critic and the advantage; the
second gives us PPO's clipping. The next three subsections are exactly those two ideas, made
precise.

\subsection{Actor and critic}

PPO trains two networks that cooperate. The \key{actor} is the policy $\pi_\theta(a\mid s)$:
given a state, it outputs a probability for each action ($\theta$ denotes its weights). The
\key{critic} is a value function $V_\phi(s)$: given a state, it estimates the return you
expect from there. The actor proposes; the critic judges whether what happened was better
or worse than expected and that judgement is what tells the actor which way to adjust.

\subsection{The advantage: was that better than expected?}

The signal the actor learns from is the \key{advantage}, how much better an action turned
out than the critic's baseline expectation. If you were in a state the critic valued at
$100$, took an action and the outcome was worth $130$, the advantage is $+30$: do that
more. If it was worth $80$, the advantage is $-20$: do that less. Estimating the advantage
well is its own small art; PPO uses \key{generalised advantage estimation} (GAE), which
blends short- and long-horizon estimates:
\begin{equation}
\hat A_t \;=\; \sum_{l=0}^{\infty} (\gamma\lambda)^{l}\,\delta_{t+l},
\qquad
\delta_t \;=\; r_t + \gamma V_\phi(s_{t+1}) - V_\phi(s_t).
\label{eq:gae}
\end{equation}
Here $\delta_t$ is a one-step surprise: the reward plus the discounted value of where you
landed, minus the value of where you were; a positive $\delta_t$ means that single step
beat expectations. GAE sums these surprises over the future, each discounted by
$\gamma\lambda$. The extra knob $\lambda$ (we use $0.98$) trades off noise against bias: near
$1$ it looks far ahead (less biased, a bit noisier), near $0$ it trusts the critic's
one-step estimate (smoother, more biased). A high $\lambda$ here gives the actor a clean,
low-bias sense of which booking decisions truly paid off many minutes later.

\paragraph{The advantage, with numbers.} Suppose the critic valued the current state at
$V(s_t) = 100$, the agent took an action, received $r_t = 5$ and the critic values the
state it landed in at $V(s_{t+1}) = 120$. The one-step surprise is
$\delta_t = 5 + 0.999\cdot 120 - 100 = 24.88$: things went markedly better than the critic
expected, so the advantage is strongly positive and PPO will make that action more likely.
Had the next state been valued at only $80$, we would get
$\delta_t = 5 + 0.999\cdot 80 - 100 = -15.08$, a negative advantage and PPO would make the
action \emph{less} likely. GAE then chains these one-step surprises over the rest of the
episode (each discounted by $\gamma\lambda$), so an action is credited not only for the
immediate surprise but for the run of good or bad surprises it set in motion. That is
exactly the right notion of credit for a booking whose benefit only unfolds across the whole
visit: the booking action inherits the advantage of every well-timed check-in it later made
possible.

\subsection{The clipped objective and why ``proximal''}

Naively, you would push the policy hard toward every action with positive advantage. The
danger is overshooting: one big greedy update can wreck a policy that took ages to shape,
and on-policy methods cannot easily undo it. PPO's whole contribution is a brake. Define
the \key{probability ratio} $\rho_t = \pi_\theta(a_t\mid s_t)\,/\,\pi_{\theta_{\text{old}}}(a_t\mid s_t)$,
how much more (or less) likely the new policy makes the action than the policy that
collected the data. PPO maximises
\begin{equation}
L^{\text{CLIP}}(\theta)\;=\;\mathbb{E}_t\Big[\min\big(\rho_t\,\hat A_t,\;\operatorname{clip}(\rho_t,\,1-\epsilon,\,1+\epsilon)\,\hat A_t\big)\Big].
\label{eq:ppoclip}
\end{equation}
Read it as a careful instruction. We want to increase probability for positive-advantage
actions, that is the $\rho_t \hat A_t$ term. But $\operatorname{clip}(\rho_t, 1-\epsilon,
1+\epsilon)$ refuses to let the ratio move outside $[1-\epsilon, 1+\epsilon]$ (with
$\epsilon = 0.2$, outside $[0.8, 1.2]$) and the $\min$ takes the less optimistic of the
clipped and unclipped versions. The effect: once an action's probability has moved about
twenty percent, the objective stops rewarding further movement on that batch, so no single
update can lurch. ``Proximal'' means each new policy stays close to the last one. This
self-imposed caution is the reason PPO is \key{stable} and stability, as the results will
show, is most of why it wins.

\paragraph{The clip, with numbers.} Say an action had a positive advantage $\hat A_t = 10$,
and the update has already raised its probability so that $\rho_t = 1.5$ (it is now $1.5$
times as likely as under the old policy). The unclipped term would reward us
$1.5\cdot 10 = 15$ for pushing even harder. But the clip caps $\rho_t$ at $1.2$, giving
$1.2\cdot 10 = 12$ and the $\min$ picks the smaller, $12$. Because the clipped value no
longer grows as $\rho_t$ rises past $1.2$, the gradient there is zero: PPO simply stops
paying us to inflate this action further on this batch. The action will keep rising, but
only a little per batch, across many batches, never in one reckless leap. For a negative
advantage the clip works symmetrically on the downside. This is the precise mechanism
behind ``small, controlled steps,'' and it is why a single noisy batch cannot wreck a
policy that took a million steps to shape.

\subsection{The training loop, in practice}

PPO does not learn from one step at a time. It runs the current policy for a fixed stretch,
a \key{rollout} of $n_{\text{steps}}$ environment steps (we use $2048$), recording every
$(s, a, r)$ and the critic's value. It then computes the advantages
(equation~\eqref{eq:gae}) for that whole batch and improves the policy on it for several
passes, $n_{\text{epochs}}$ (we use $8$ to $10$), chopping the rollout into mini-batches of
a few hundred. Then it throws the data away and collects a fresh rollout with the improved
policy. This rollout-then-refine rhythm is what ``on-policy'' looks like in practice: the
data always reflects the current policy, which is why it cannot be hoarded the way DQN
hoards its replay buffer. One more ingredient keeps exploration alive: an \key{entropy
bonus} (a small reward, weight $\approx 0.01$, for keeping the action probabilities spread
out rather than collapsing to a single certain choice too early). Without it a policy can
prematurely commit to a mediocre habit; with it, the agent keeps a little curiosity until
the evidence is in. The far-sighted discount ($\gamma = 0.999$ for booking), the low-bias
advantage ($\lambda = 0.98$), the clip ($\epsilon = 0.2$) and this entropy nudge are the
handful of dials that define our PPO; none of them is exotic and the defaults work because
masking has already made the problem clean.

\subsection{How PPO differs from Q-learning}

Two differences are worth stating plainly. First, PPO is \key{on-policy}: it learns from
data collected by the current policy, then throws it away, whereas Q-learning and DQN are
\key{off-policy} and reuse old experience. On-policy data is more expensive but better
matched to what the policy is currently doing. Second, PPO adjusts probabilities smoothly
and never takes a $\max$ over noisy value estimates, so the overestimation curse of
section~\ref{sec:doubleq} simply does not arise for it. That single fact will explain a lot about why PPO
is steady where DQN is jumpy.

% =============================================================================
\section{Telling the agent the rules of the building: action masking}
% =============================================================================

The museum is a graph: from any room only a few moves are physically possible, but the
action head of the network has about a hundred outputs, one per possible room. Without
help, the policy spends most of its probability on moves that do not exist and has to
\emph{learn}, from penalties, that they are illegal, wasting effort on the obvious.

\key{Action masking} hands the agent the rulebook for free. Before turning the network's
raw scores (logits) into probabilities, we set the score of every illegal action to minus
infinity, so it receives exactly zero probability:
\begin{equation}
\pi_\theta(a\mid s)\;=\;\frac{\exp(z_a)\;\mathbb{1}[a\in\mathcal A(s)]}{\sum_{a'\in\mathcal A(s)}\exp(z_{a'})},
\label{eq:mask}
\end{equation}
where $z_a$ is the network's score for action $a$, $\mathcal A(s)$ is the set of legal
actions in state $s$ and $\mathbb{1}[\cdot]$ is $1$ for legal actions and $0$ otherwise.
The denominator sums only over legal actions, so the probabilities of the legal moves
renormalise to one and every illegal move is impossible by construction. \key{MaskablePPO}
is just PPO with this mask applied identically during training and evaluation.

This sounds like a minor convenience. It is in fact the single most important
architectural choice in the project, because it removes an entire learning problem (``which
moves are even allowed?'') and lets every scrap of gradient go into the decision that
actually matters.

\paragraph{Why the unmasked version is genuinely harder, not just messier.} Without
masking, the network must \emph{learn} to avoid illegal moves from a penalty and that has
three costs that compound. It wastes capacity: probability mass and gradient go into
suppressing the eighty-odd impossible actions instead of ranking the handful of real ones.
It slows learning: early on, most sampled actions are illegal, so most steps teach the dull
lesson ``that was not allowed'' rather than anything about good play. And it adds noise to
the advantage estimates, because an episode's return now mixes the signal of good decisions
with the static of accidental illegal moves and their penalties. Masking is, mathematically,
exactly equivalent to handing the agent a different, smaller, state-dependent action space,
the legal one, so none of these costs is ever paid. That is why we describe it not as a
tweak but as changing the problem the network has to solve. The results section quantifies
exactly how much it buys and, revealingly, \emph{where}: only where the underlying decision
is hardest.

% =============================================================================
\section{The value-based alternative: DQN}
% =============================================================================

For contrast, the project also trains \key{DQN} (Deep Q-Network), the neural version of
Q-learning. It keeps the value-based spirit: learn $Q(s,a)$ with a network $Q_\theta$, then
act greedily. It minimises the squared temporal-difference error, the same target as
equation~\eqref{eq:qupdate} but fit by gradient descent:
\begin{equation}
L(\theta)\;=\;\mathbb{E}\Big[\big(\,r+\gamma\max_{a'}Q_{\theta^-}(s',a')-Q_\theta(s,a)\,\big)^2\Big].
\label{eq:dqn}
\end{equation}
Two tricks make this stable enough to work with a network. A \key{replay buffer} stores
past experience and replays random batches of it, so the network is not whiplashed by the
latest correlated steps. And a \key{target network} $Q_{\theta^-}$, a slowly updated copy
of the weights, supplies the value of $s'$ in the target, so the thing we are chasing does
not move every single step. We give DQN a buffer of a hundred thousand to two hundred
thousand transitions, a slow target update and a sensible exploration schedule.

\paragraph{The deadly triad.} There is a well-known warning in reinforcement learning that
three ingredients, when combined, make learning unstable: \key{bootstrapping} (updating an
estimate toward a target that contains another estimate), \key{function approximation} (a
neural network rather than an exact table) and \key{off-policy} learning (training on data
from a different, older policy via the replay buffer). This trio is called the \key{deadly
triad} and DQN has all three at once. None of them is individually fatal, the tabular toy
had bootstrapping and off-policy learning and was fine, but together, on a hard noisy
problem, they can let errors feed on themselves: a wrong value gets bootstrapped into
neighbouring states, the network generalises the error to states it has not even seen and
the stale buffer keeps reinforcing it. PPO sidesteps the triad entirely: it is on-policy and
does not bootstrap a $\max$, which is a large part of why it is the calmer learner.

\paragraph{The overestimation, returning.} Notice the $\max_{a'}$ sitting right there in
equation~\eqref{eq:dqn}. This is the same maximisation that caused overestimation in
section~\ref{sec:doubleq} and plain DQN does \emph{not} use the double-table fix (its deep cousin
\key{Double-DQN} exists precisely to add one, by choosing the next action with the online
network but scoring it with the target network, exactly the Double-Q trick of section~\ref{sec:doubleq-fix}).
On a deterministic toy this did not matter, which is why Double-Q tied vanilla there. On the
full museum, where the crowd is genuinely stochastic and the value targets are therefore
noisy, it matters a great deal. Our DQN is also unmasked and chases a moving target with a
big function approximator. Keep all of this in mind: bootstrapped overestimation, the deadly
triad, no masking. It is precisely the recipe for the crowd-fragile, high-variance behaviour
we will see in the numbers and it is the live, full-scale echo of the toy lesson that
Double-Q only earns its keep under noise.

% =============================================================================
\section{Exploration versus exploitation, the running tension}
% =============================================================================

Every method in this guide faces the same dilemma at every step. Should the agent take the
move that looks best on current evidence (\key{exploit}), or try a move it is unsure about to
find out whether it is secretly better (\key{explore})? Pure exploitation locks onto the
first decent habit and never improves, because it never gathers evidence about the roads not
taken. Pure exploration learns a lot but never cashes any of it in. Good learning explores
heavily early, when it knows little and shifts toward exploiting as the evidence piles up.

The tabular and value-based methods (Q-learning, DQN) manage this with
\key{$\varepsilon$-greedy}: with probability $\varepsilon$ act at random, otherwise take the
best-known move. We start $\varepsilon$ near $1$ (almost all exploration) and decay it toward
near $0$ (almost all exploitation), so the agent is adventurous as a beginner and disciplined
as it matures. PPO reaches the same goal by a different door, the \key{entropy bonus}: a small
reward for keeping the action probabilities spread out rather than collapsing onto one
near-certain choice. Entropy is just a measure of how undecided a distribution is, high when
the options are balanced, low when one dominates, so paying for entropy is literally paying
the policy to stay a little open-minded until the evidence justifies committing.

These are two costumes for one idea and getting it wrong is behind real failures. The
do-nothing trap (the free decline) is exactly an exploration failure: the safe action paid off
immediately, the agent stopped trying the booking alternative before it ever experienced the
delayed reward and so it never learned that booking was better. Whenever an agent ``gets
stuck'' on a mediocre habit, exploration is the first thing to interrogate.

% =============================================================================
\section{The three deep methods at a glance}
% =============================================================================

Before reading the results, it helps to line up the three networked methods on the axes that
actually drive their behaviour. Everything in the table traces back to the sections above,
and every row will explain a feature of the numbers.

\renewcommand{\arraystretch}{1.35}
\begin{center}
\begin{tabular}{p{3.2cm} p{3.5cm} p{3.5cm} p{3.5cm}}
\toprule
 & \textbf{DQN} & \textbf{PPO} & \textbf{MaskablePPO} \\
\midrule
Family & value-based & policy-gradient & policy-gradient \\
Learns & $Q(s,a)$, acts greedily & the policy directly & the policy directly \\
On- or off-policy & off-policy (replay) & on-policy & on-policy \\
Takes a max over noisy values? & yes (the liability) & no & no \\
Knows legal moves for free? & no & no & yes (masking) \\
Main strength here & sample-reuse & stable updates & stable + clean action set \\
Main weakness here & overestimation under noise & must learn legality & none of note \\
\bottomrule
\end{tabular}
\end{center}

Read the table as a set of predictions. The ``yes'' in the max-over-noisy-values row is why
DQN alone suffers overestimation on the noisy crowd. The ``no'' in the legal-moves row for
plain PPO is why it wastes effort and stumbles exactly where the decision is hardest. And the
clean column for MaskablePPO, stable updates and a free-of-charge legal action set, is why it
wins. The results section is essentially this table, made quantitative.

% =============================================================================
\section{The scaffolding: curriculum, common random numbers and seeds}
% =============================================================================

Three pieces of machinery are not algorithms but make the comparisons trustworthy and the
project's conclusions lean on them.

\paragraph{Curriculum.} Training an agent directly against the full five-thousand-person
crowd is so noisy early on that it learns slowly. So for the navigation agent we ramp the
difficulty: a first stage at about a third of the peak crowd, then two-thirds, then the
full crowd, always over a complete museum day. Learning the easy version first and then
transferring is faster and more reliable, the same reason you would not teach someone to
drive in rush-hour traffic on day one. The deeper justification is about signal-to-noise. A
beginner policy makes mostly poor choices and in a packed museum those poor choices are
buried under heavy crowd randomness, so the learning signal is weak and the agent flails.
In a light crowd the same poor choices produce clearer, less noisy consequences, so the
basics (where the masterpieces are, how to pace, when to leave) lock in quickly; once they
are in place, the agent can withstand the full crowd's noise and refine. The ordering
matters because skills learned early are a scaffold for the harder ones, not a separate
lesson to be unlearned. Note that the booking agents, whose decision is simpler than full
navigation, do not need the ramp and are trained directly at each crowd level.

\paragraph{Common random numbers.} To compare two policies fairly, they must face the
\emph{same} crowds. We fix six crowd scenarios (seeds $900000$ to $900005$) and score every
policy on those identical six. This is \key{common random numbers}: because both policies
meet the same randomness, any difference in their scores reflects the \key{policy}, not the
luck of which crowd each happened to draw. It sharply reduces the noise in a comparison.

\paragraph{Why it helps, concretely.} Imagine policy A scores $3700$ on one day and policy B
scores $3600$ on another day. Is A genuinely better, or did it just draw an easier crowd? You
cannot tell; the difference could be all luck. Now force both to play the \emph{same} six
crowd scenarios and compare them scenario by scenario: if A beats B on the very same crowd,
the gap is about the policies, because the only thing that changed is the policy. It is the
logic of a fair race, everyone runs the same track on the same day, so the finishing order
means something. This is why our evaluation crowds (seeds $900000$ through $900005$) are
frozen and identical for every model in the grid and it is what lets us read a small gap
between two algorithms as real rather than noise.

\paragraph{Multiple seeds.} A single training run is one sample of a random process (random
network initialisation, random exploration, random mini-batches). One run can be lucky or
unlucky, so a single number can mislead. We therefore train each cell of the comparison
under five different \key{training seeds} ($1, 3, 7, 42, 202606$) and report the mean and
the standard deviation across them. The mean is our best estimate; the standard deviation
tells us how reliable that estimate is. As we will see, this matters: one cell of the grid
looks like a clean win on a single seed but is revealed, across five, to be a coin-flip
between a large win and a collapse.

\paragraph{How to read a training curve.} Several figures in the project plot a learning
curve and they all read the same way. The horizontal axis is training progress (episodes
for the toy, environment steps for the deep agents); the vertical axis is the return the
policy achieves. A healthy curve starts low (the untrained agent is bad), climbs as the loss
is nudged downhill batch after batch and then \key{plateaus} once the policy stops
improving, that flat top is \key{convergence} and it is the part we care about, not the
noisy beginning. The shaded band around the line is variability (a rolling confidence
interval, or the spread across seeds): a narrow band means a reliable result, a wide one
means a fragile one. Two warnings the curves teach. First, read the plateau, not the single
highest point; the peak of a noisy curve is a lucky fluctuation, not the policy's true level.
Second, a curve that climbs nicely in \emph{training} reward can still hide a bad
\emph{evaluation} policy if training and evaluation differ (the very issue behind one of our
failures, where a stochastic policy trained well but its greedy version evaluated poorly).
A learning curve is a thermometer, not a verdict; the verdict comes from the common-random-
number evaluation above.

% =============================================================================
\section{The decision that actually matters: the booking problem}
% =============================================================================

Early in the project the agent was asked to learn how to \emph{walk} the museum. That was
a mistake we corrected and the correction is itself a lesson worth stating: a real
visitor has a map and is guided by staff, so optimising the walk is a shortest-path puzzle
dressed up as learning, with a trivial, fully observed state. The genuinely uncertain,
consequential decision is the \key{booking} and reframing the task around it is what made
the project's centerpiece. The walking then happens underneath as automatic pacing.

\paragraph{Choosing the MDP is the modelling.} It is tempting to think the hard part of a
reinforcement-learning project is the algorithm. It is not. The hard part is deciding
\emph{what decision the agent is even making}, that is, choosing the state, the actions and
the reward. We spent months making a navigation agent work and the brittle result kept
whispering that something was wrong. The realisation was that we had chosen the wrong MDP:
walking a known map with a guide is barely a decision at all, so no algorithm, however good,
could make it interesting. Moving the action from ``which room next'' to ``when and how far
ahead, to book'' turned a trivial shortest path into a genuine sequential decision under
uncertainty with a delayed payoff, the kind of problem where learning earns its keep. The same equations and the same PPO then produced a result worth reporting,
not because the algorithm got better but because the question did. If you take one
methodological lesson from this guide, let it be that the formulation is more than half the
work: a beautiful optimiser pointed at the wrong decision optimises nothing.

\subsection{A two-phase episode}

Each episode now has two phases. In \key{Phase 1 (booking)}, before the visit, the agent
chooses a \key{lead time}, how many days in advance to book, from $\{1, 7, 21, 35\}$ and
then, for each of the three gated masterpieces, whether to take the offered slot. In
\key{Phase 2 (the visit)}, it walks the museum and must \key{check in} at each painting
inside the time window it booked. The discount is set high, $\gamma = 0.999$, because the
reward for a good booking does not arrive until check-in, which can be hundreds of minutes
into the visit; a far-sighted agent is mandatory.

\paragraph{One full episode, narrated.} Picture a packed day. In Phase 1 the agent, reading
that the museum is full, chooses a lead time of $35$ days, because the fill curve tells it
that anything shorter will not secure the popular slots. It accepts the offered windows for
Botticelli, Leonardo and Raphael, spread across the day. So far it has earned almost
nothing; the payoff is still hundreds of steps away. Phase 2 then begins at the entrance. The
agent walks room by room (every move a legal neighbour), pacing itself so that it arrives at
Botticelli inside its ten-minute window, where it finally banks the appreciation and a
check-in bonus; it continues to Leonardo and Raphael in their windows, drops to the first
floor for Caravaggio and leaves before closing. Only now, at each check-in, does the large
reward land and the high discount carries that reward all the way back to justify the
lead-$35$ choice made before the visit even started. That backward flow of credit, from a
check-in late in the afternoon to a booking made weeks in advance, \emph{is} the learning
problem in miniature and it is why the next subsection is about credit assignment.

\subsection{The reward, component by component}

The per-minute reward is a sum of terms and reading them is the fastest way to understand
what the agent is being asked to want:
\begin{equation}
r \;=\; \underbrace{r_{\text{art}}}_{\text{appreciation}} \;+\; \underbrace{r_{\text{closing}}}_{\text{get-out pressure}} \;+\; \underbrace{r_{\text{completion}}}_{\text{finished a must-see}} \;+\; \underbrace{r_{\text{checkin}} + r_{\text{wait}}}_{\text{booking discipline}} \;+\; \underbrace{r_{\text{time}}}_{\text{small step cost}} .
\label{eq:reward}
\end{equation}
The central term is the appreciation,
\begin{equation}
r_{\text{art}} \;=\; \text{importance} \,\times\, \underbrace{\frac{1}{1+\alpha\,\text{density}^2}}_{\text{crowd discount}} \,\times\, \underbrace{\text{novelty}}_{\text{satiation}} ,
\label{eq:rart}
\end{equation}
which says a room pays its \key{importance}, scaled down when it is crowded (the crowd
discount, with $\alpha$ small so a crowded masterpiece still beats skipping it, $\alpha=0.5$
for the art lover, $1.0$ for the tourist) and scaled down the longer you have already
stood there (\key{novelty} decays, so value \key{satiates} and the agent eventually moves
on).

\paragraph{The crowd discount, with numbers.} Take a masterpiece of importance $10$ and the
art lover's crowd-aversion $\alpha = 0.5$. In an empty room (density $0$) the appreciation
is $10 \times \tfrac{1}{1+0} = 10$ times novelty. In a room jammed to twice capacity
(density $2$) it is $10 \times \tfrac{1}{1 + 0.5\cdot 2^2} = 10\times\tfrac{1}{3} \approx 3.3$
times novelty: the same painting pays a third as much when packed. The number that matters
is that it is \emph{still positive}, so a crowded Botticelli still beats skipping it. This is
a deliberate calibration. An earlier, harsher crowd term drove the discounted value below the
reward of simply walking past and then the optimiser rationally skipped the masterpiece,
producing a ``fake art lover'' who speed-ran the famous rooms. Shrinking $\alpha$ until a
crowded masterpiece is still worth entering is what made the modelled connoisseur behave like
one. As for novelty: it decays with each minute you stay, so after enough time in a room the
marginal minute pays almost nothing and moving on becomes the better action without any
explicit ``you are bored now'' rule. Satiation, not a timer, is what ends a visit to a room. The other terms each fix a specific failure: $r_{\text{closing}}$ is a pressure that
grows as closing nears and you are far from an exit, so the agent leaves on time;
$r_{\text{completion}}$ is a one-off bonus for fully appreciating a must-see;
$r_{\text{checkin}}$ rewards arriving inside a booked window and $r_{\text{wait}}$ gently
penalises idling for one; $r_{\text{time}}$ is a small per-step cost so dithering is never
free. The booking phase adds a few more: a penalty for declining a masterpiece (since these
profiles always want all three), a no-show penalty for booking a slot and missing it and,
decisively, an \key{immediate off-pace penalty} $\text{mismatch}_k \cdot |\text{window} -
\text{estimate}|$ paid at booking time. That last one matters for a reason we come to
below.

\subsection{The fill curve and why ``book early'' is a real decision}

Booking is only interesting because slots are scarce on busy days. A slot is available only
if the lead time is long enough:
\begin{equation}
\text{lead} \;\ge\; \text{lead}_{\max} \times \text{popularity} \times \frac{\text{daily visitors}}{2500}.
\label{eq:fill}
\end{equation}
The busier the day, the larger the right-hand side, so the longer ahead you must book to
secure a popular masterpiece. Book late on a packed day and the good slots are simply gone,
and you miss the painting entirely, which is a large loss of reward. Book early and you
lock in well-spaced windows. There is no downside to booking early; the only question is
how early and the agent has to \emph{discover} that the answer rises with the crowd. This
is the heart of the centerpiece result.

\subsection{Why this is hard: the credit-assignment problem}

Booking is a textbook case of the hardest thing in reinforcement learning, \key{credit
assignment} across a long delay. The agent picks a lead time in the first few steps of the
episode. The consequence, did it actually get into the masterpiece inside a good window, or
did it arrive to find the slot mistimed and waste the visit, does not reveal itself until
hundreds of minutes later at check-in. So the reward that should teach ``you booked too
late'' is separated from the booking action by a vast stretch of unrelated walking. A
far-sighted discount ($\gamma = 0.999$) is necessary to carry the signal back that far, but
it is not sufficient: a faint signal smeared across four hundred steps produces a tiny,
noisy gradient on the booking action and early versions of the agent only half-learned the
lesson (a blurry mix of lead times, two-and-a-bit masterpieces secured at the packed crowd).
The cure was to manufacture a \key{local proxy} for the distant outcome: the immediate
off-pace penalty $\text{mismatch}_k \cdot |\text{window} - \text{estimate}|$, paid right at
booking time, which is large exactly when the booked window is far from when the agent can
realistically arrive. Now the agent feels a consequence \emph{immediately}, correlated with
the distant one and the gradient has something nearby and strong to follow. This single
trick is what turned ``book early'' from half-learned to crisply learned and it is a
beautiful illustration of a general principle: when the true reward is too far away to
teach, engineer a faithful local signal that points the same way.

\subsection{Two modelling choices worth being explicit about}\label{sec:modelling}

Two design decisions shape these results and deserve to be named, because a careful reader
should know which parts are learned and which are assumed.

First, the agent is \emph{not} allowed to decline all three masterpieces. We tried leaving
a free ``decline'' button and both visitor types collapsed to declining everything (we
explain why in the failures section: a free skip is a classic do-nothing trap). We treat
``always books the three masterpieces'' as a \key{type characterisation}, a fact about who
these visitors are, nobody flies to Florence and skips all of Botticelli, Leonardo and
Raphael, rather than something the agent must rediscover. The genuinely learned decision is
the \emph{lead time} and that is what we report.

Second, the agent needs an estimate of when it will arrive at each masterpiece in order to
judge a booking and that estimate is computed once from reference walks and then
\key{frozen} during training. The reason is subtle and important: if the arrival estimate
were updated online from the agent's own past arrivals, it would become a function of the
policy that is itself feeding the policy, a self-referential loop that drifts (again, see
the failures section). Freezing it breaks the loop. This is an assumption, but a transparent
one and the qualitative lesson (book earlier when busier) does not depend on its exact
value.

% =============================================================================
\section{Reading the results through the theory}
% =============================================================================

This is the section the rest of the guide was building toward. For each finding, we say
what the number is and then \why{why the theory above forces it to be that way.} If you
understand these, you can look at any of the project's tables and explain it.

\subsection{Why Q-learning and Double-Q tie on the toy (83.3 vs 83.1)}

\why{Because overestimation, the only thing Double-Q fixes, is born of noise and the toy
has none.} The crowd waves are deterministic and the rewards are a fixed function of
(state, action), so at convergence there is no random jitter in the TD target for the
$\max$ in equation~\eqref{eq:qupdate} to inflate (equation~\eqref{eq:jensen} bites only when
the estimates are noisy). Double-Q's two-table machinery therefore cancels a bias that is
not present, changing nothing on average. The tiny gap between runs is exploration noise,
which both share equally. Predicted consequence: the moment we move to a genuinely
stochastic environment, this tie should break in a way that hurts the single-table method,
and it does, in DQN below.

\subsection{Why the handcrafted baselines lose and the ``smart'' ones lose worst}

The learned policies score around $83$; the best hand rule, \key{random}, manages only
about $35$ and the structured rules are \emph{negative}: \key{peak\_avoidance} $\approx -34$,
\key{greedy\_least\_crowded} $\approx -43$, \key{default\_path} $\approx -61$,
\key{greedy\_value\_ratio} $\approx -72$. \why{The negatives come from the closing-time and
step-cost terms in the reward: a fixed rule that mistimes its visit gets caught wandering
as closing nears and pays the egress pressure, while collecting little appreciation because
it stood in crowds.} The reason the \emph{structured} rules do worse than \emph{random} is
the sharpest point: \key{greedy\_value\_ratio} marches straight at the highest-value room
\emph{regardless of when}, so it reliably arrives at Botticelli at the peak, exactly the
worst moment and the crowd discount in equation~\eqref{eq:rart} guts its reward. A rule
that is confidently wrong about timing is worse than one that stumbles around, because the
random walker at least samples some off-peak minutes. This is the toy's entire moral:
\emph{when} you visit dominates \emph{whether} and only the learner discovers it.

\subsection{Why MaskablePPO wins everywhere}

MaskablePPO is best or tied-best in every cell of the grid. \why{Two of its properties from
the sections above combine. First, masking (equation~\eqref{eq:mask}) deletes the entire
sub-problem of learning which moves are legal, so all of its capacity goes into the booking
and pacing decision. Second, the clipped objective (equation~\eqref{eq:ppoclip}) makes its
updates small and non-destructive, so it converges to a good policy and stays there rather
than thrashing.} It has no $\max$ over noisy values, so it never suffers the overestimation
that destabilises DQN. Steady optimisation of the right, constrained problem is simply hard
to beat.

\subsection{Why unmasked PPO matches at easy crowds but collapses at the packed tourist}

This is the cleanest demonstration of what masking buys. Plain PPO equals MaskablePPO at
the quiet and busy tourist ($\approx 3686$ and $3695$), then at the packed tourist it
collapses to $2093$, \emph{below} the no-booking baseline of $2680$, while MaskablePPO holds
$3685$. \why{At low crowd the booking problem is forgiving: many lead times work, so masked
and unmasked policies converge to the same easy optimum. At the packed crowd the fill curve
(equation~\eqref{eq:fill}) is at its tightest, only a long lead and precise pacing secure
the slots and that knife-edge is exactly where wasting probability on illegal or pointless
moves is fatal. The unmasked policy cannot reliably hold the tight requirement and degrades
into something worse than not booking at all.} The masking is decisive precisely where the
decision is hard, which is why the gap appears only in that corner of the grid and is
reproducible across all five seeds.

\subsection{Why DQN is crowd-fragile and high-variance}

DQN's value-based baseline on the art lover falls from about $4527$ at the quiet crowd to
$2044$ at the packed crowd and its bars carry the widest error bars in the figure, with
standard deviations up to roughly $\pm 2600$ across seeds. \why{This is section~\ref{sec:doubleq}'s
overestimation returning to collect its debt. The full museum is genuinely stochastic, so
the value targets are noisy; the $\max$ in equation~\eqref{eq:dqn} then inflates them
(Jensen again) and plain DQN, unlike Double-Q, has no mechanism to cancel the bias. The
inflation grows with the crowd because a busier museum is a noisier one, which is why the
collapse tracks the crowd. And because value-based bootstrapping off a moving target is
sensitive to initialisation and replay order, different seeds land in very different
places, hence the large spread.} Everything that made DQN look fine on the deterministic
toy is exactly what makes it fragile here; the contrast is the point of including it.

\subsection{Why ``book early when busier'' emerges from learning}

The learned lead time rises with the crowd: the art lover goes from booking $7$ days ahead
when quiet to $35$ days ahead when busy; the tourist from $7$ to about $21$. \why{The fill
curve makes late booking on a packed day forfeit the masterpieces, a large reward loss
under equation~\eqref{eq:rart} and the far-sighted discount $\gamma = 0.999$ lets that
distant loss reach back to the booking action. But a delayed loss alone is a weak teacher,
the gradient for a payoff hundreds of steps away is faint, which is why early versions only
half-learned it. The immediate off-pace penalty $\text{mismatch}_k\cdot|\text{window}-
\text{estimate}|$ is what fixed it: it puts a small, local cost at booking time that
correlates with the distant outcome, giving the gradient something nearby to follow.} The
behaviour is not hand-coded; it is the optimum of the reward, made findable by one
well-placed local signal.

\subsection{Why the tourist is rock-stable but the art lover at max crowd is a coin-flip}

This is the subtlest reading and the one the five seeds exist to reveal. The tourist's
intervened scores have essentially zero spread across seeds at every crowd and its lift is
a steady $+36\%$ to $+38\%$. The art lover's lift is healthy at the quiet and busy crowds
($+39\%$, $+29\%$) but at the packed crowd it averages only $+20\%$ \emph{with a huge
standard deviation} of about $\pm 1900$: four of the five seeds reach roughly $+41\%$ and
one seed collapses to a score well below the baseline.

\why{The difference is the shape of the optimisation landscape for each profile. The tourist
has short dwell times, cares about only a handful of recognisable rooms and is only mildly
crowd-averse, so its optimal visit is simple and essentially unique; every seed rolls down
to the same solution and with deterministic evaluation that yields identical scores.} The
art lover is the opposite: long dwells, value in many rooms and strong crowd-aversion, so
at the packed crowd the optimum is a knife-edge, book the longest lead \emph{and} pace a
long, crowd-sensitive visit so every window is met before closing. \why{That narrow optimum
sits next to a basin of failure and which one a given training run falls into depends on
its seed. Four seeds find the knife-edge ($+41\%$); one slips into the failure basin and
collapses. Averaging gives the modest $+20\%$ with the wide band.} The right way to report
this cell and the way the figure now shows it, is the mean with its error bar: the
intervention usually delivers a large gain for the art lover even at maximum crowd, but at
that single hardest cell the result is not guaranteed and the spread says so.

\subsection{Why all three masterpieces are secured at every crowd}

Across every cell the agent ends up with $3/3$ masterpieces booked and checked in.
\why{This is the do-nothing trap, resolved by design (section~\ref{sec:modelling}): because declining is
not an available action, the only question the agent optimises is the lead time and with
the fill curve the way to keep all three is simply to book early enough.} So securing all
three is not a surprising feat of intelligence; it is the consequence of removing an
unrealistic escape hatch and letting the agent solve the one real problem, when, rather than
a fake problem, whether.

\subsection{Why the learned routes are clean (zero teleports)}

When we audit the walks, every room-to-room move is a physical one-hop neighbour; there are
no jumps across the map. \why{This is action masking doing its quiet job: the mask makes a
non-adjacent move literally impossible (zero probability, equation~\eqref{eq:mask}), so a
teleport cannot occur even by accident.} A reader who did not know about masking might be
impressed that the agent ``learned the floor plan''; in truth it was handed the floor plan,
which is exactly the point of masking and exactly why we do not have to spend the agent's
capacity teaching it geography.

\subsection{Why lead times jump in discrete steps}

The learned lead time does not glide smoothly; it sits at $7$, then jumps to $35$, with
little in between. \why{That is because the action set offers only $\{1, 7, 21, 35\}$ days,
and the fill curve (equation~\eqref{eq:fill}) is a threshold: below the required lead the
slot is simply unavailable, above it the slot is secured. So the agent has no reason to pick
a middle value, it picks the smallest offered lead that clears the threshold for the current
crowd.} The discreteness of the curve in the figures is therefore a faithful picture of a
discrete decision, not a plotting artifact and the jump to $35$ at the packed crowd is the
agent reporting that nothing shorter would clear the threshold.

% =============================================================================
\section{One-page cheat-sheet}
% =============================================================================

\renewcommand{\arraystretch}{1.35}
\begin{center}
\begin{tabular}{p{2.7cm} p{5.2cm} p{6.4cm}}
\toprule
\textbf{Method} & \textbf{What it is, in one line} & \textbf{Why its number looks the way it does} \\
\midrule
Tabular Q-learning & Fill a table of long-run move values by trial-and-error nudges. & Beats every hand rule ($\approx 83$) because it learns to visit off-peak; timing dominates value. \\
Double Q-learning & Q-learning with two tables to cancel optimism. & Ties vanilla ($\approx 83$) on the toy: deterministic, so no overestimation to remove. \\
Handcrafted rules & Fixed heuristics, no learning. & Mostly negative; the confidently mistimed ones (greedy value) lose worst, below random. \\
PPO & Adjust the policy directly with small, clipped steps. & Stable; no noisy $\max$, so it never overestimates. The backbone of the wins. \\
MaskablePPO & PPO told which moves are legal before choosing. & Best or tied everywhere; masking frees all effort for the real decision. \\
DQN & Neural Q-learning with replay and a target network. & Crowd-fragile, high-variance: the overestimation that was harmless on the toy bites under a noisy crowd. \\
Action masking & Zero out illegal actions before the softmax. & Decisive only at the packed tourist, where the booking constraint is tightest. \\
Curriculum / CRN / seeds & Ramp difficulty; fix the test crowds; repeat runs. & Make the comparison fair and reveal the art-lover/max-crowd cell as a coin-flip, not a clean win. \\
\bottomrule
\end{tabular}
\end{center}

\vspace{1em}
\noindent\emph{The single thread through all of it:} value-based methods take a maximum
over their own noisy estimates and so flatter themselves when the world is noisy;
policy-gradient methods nudge probabilities and do not. That one difference, plus the free
gift of action masking, explains the ranking of the algorithms and the shape of each
profile's reward landscape explains the spread of the numbers within the winner.

% =============================================================================
\section{Failures as reinforcement-learning lessons}
% =============================================================================

Most of the project was spent being wrong in instructive ways and each mistake maps onto a
named, general phenomenon in reinforcement learning. Understanding these is as valuable as
understanding the algorithms, because they are the traps you fall into whenever you try to
make an agent want the right thing. We present five, each as a symptom, the underlying
concept and the fix.

\subsection{The unlearnable reward (sparse, all-or-nothing signals)}

\emph{Symptom.} The agent refused to leave the museum; it would squat in a room until the
day ended. The obvious fix, ``zero out all your reward if you fail to exit,'' made it
\emph{worse}.

\emph{Concept.} A reward that is all-or-nothing at the very end is \key{sparse} and, worse,
nearly flat: almost every policy fails to exit, so almost every policy gets the same zero,
and a function that is flat across policies has \emph{no gradient}. Gradient descent needs a
slope to walk down; an all-or-nothing terminal wipe gives it a cliff surrounded by a
plateau and from the plateau nothing points toward the cliff edge. The agent cannot learn
what it cannot feel a gradient toward.

\emph{Fix.} Replace the wipe with a \key{dense} signal: a per-step egress pressure that grows
smoothly as closing nears and as you get further from an exit, plus a finite (not infinite)
penalty for not leaving and a bonus for leaving. Now every step toward the door is rewarded a
little, so there is a continuous slope leading out. The lesson, which recurs throughout the
field, is that a reward must be \emph{shaped} so that better behaviour is always at least
slightly better-scoring; sparse rewards are correct but often unlearnable.

\subsection{Acting on what you cannot see (partial observability)}

\emph{Symptom.} The rule ``stay in a room until you have appreciated it enough, then move
on'' never stabilised; the agent over- and under-stayed seemingly at random, even with a
sensible reward.

\emph{Concept.} The right action depended on \emph{how much value the agent had already
extracted from the current room}, but that quantity was not in the observation. The decision
was therefore not \key{Markov} in the state the policy could see (recall section~\ref{sec:vocab}): the
agent was being asked to act on a hidden variable. This is a \key{partially observed} problem
(a POMDP) and no amount of training fixes it, because the information simply is not there to
condition on.

\emph{Fix.} Put the missing variable into the state. We added a per-room
\key{appreciation-progress} vector to the observation, so ``have I seen enough of this room?''
became a function of something the policy can actually read. The lesson is the most basic one
in the vocabulary section, made painful by experience: the state must contain everything the
optimal action depends on, or you are not solving the MDP you think you are.

\subsection{Reward shaping that backfired (reward hacking)}

\emph{Symptom.} Well-meant extra reward terms made behaviour worse. A ``boredom'' penalty
meant to discourage lingering produced \emph{looping}; a strong pull toward completing a tour
produced \emph{ping-pong} between two rooms.

\emph{Concept.} Every term you add to a reward is a new objective and the optimiser will
satisfy it in whatever way scores highest, including ways you did not intend. This is
\key{reward hacking}: circling a loop technically avoided the boredom penalty more cheaply
than exiting did; an over-weighted completion pull made oscillating between two ``progress''
rooms score better than actually finishing. The agent was not malfunctioning; it was
obeying, too literally, a reward we wrote carelessly.

\emph{Fix.} Remove the boredom penalty entirely and reduce the tour pull to a gentle nudge,
letting the satiation term carry ``stop lingering'' on its own. The lesson: prefer the
\emph{minimal} shaping that points the right way and stress-test every term by asking
``what is the cheapest way for a ruthless optimiser to collect this and is that what I
want?''

\subsection{The do-nothing trap (a safe action with delayed reward)}

\emph{Symptom.} Given a free ``decline this masterpiece'' button, both visitor types learned
to decline all three, even though booking was worth far more and even after a curriculum
that first showed booking paying off.

\emph{Concept.} Declining is \emph{immediately safe}, costs nothing right now, while the
reward for booking is \emph{delayed} to check-in. Faced with a safe-now option versus a
better-but-later one, an optimiser that has not yet learned the delayed payoff will park on
the safe action and once parked it stops exploring the alternative, so it never discovers
the payoff. This is a \key{do-nothing local optimum}, a close relative of the
exploration problem: the agent is stuck on a hill that is locally fine and globally poor.

\emph{Fix.} Here we judged the action itself unrealistic for these visitor types and removed
it from the formulation, treating ``always books'' as a type trait (section~\ref{sec:modelling}). The
general lesson is that a voluntary action that is always immediately safe but only
delayed-rewarding will be abused unless either the formulation rules it out or exploration is
strong enough to bridge the delay.

\subsection{The self-referential runaway (feedback into your own state)}

\emph{Symptom.} The agent's estimate of its own arrival time drifted later and later each
update, until it ``planned'' to arrive at closing, even on a fixed, sensible walk.

\emph{Concept.} The arrival estimate was an online average of the agent's \emph{own} past
arrivals and it fed back into the policy that produced those arrivals. A quantity that is a
function of the policy's history, feeding the policy, is a feedback loop; any small drift
compounds, like an alarm clock you keep resetting to whenever you happened to wake up. It
also makes the environment \key{non-stationary} from the agent's point of view: the rules it
is learning against are quietly changing because of its own behaviour.

\emph{Fix.} \key{Freeze} the prior: compute the arrival estimate once from reference walks and
hold it fixed during training. Breaking the loop restores a stationary target. The lesson:
beware any state variable that is a function of the policy's own past and feeds back into the
policy; either freeze it or model the feedback explicitly.

% =============================================================================
\section{Common confusions, answered}
% =============================================================================

A handful of questions come up every time these results are presented. Answering them
collects the guide's ideas into a few sharp points.

\paragraph{Is a bigger reward number always the goal?} No. The reward is a tool we
\emph{designed} to point the agent at behaviour we want and the agent will maximise exactly
what we wrote, loopholes included (the reward-hacking failure). A higher number is only
``better'' if the reward faithfully encodes good behaviour. Half the project was getting the
reward to mean what we intended, not getting the agent to maximise it.

\paragraph{Why not just compute the optimal visit with a solver?} Because it is not a static
puzzle. The crowd is stochastic and only partially observed, the payoff for a booking is
delayed by hundreds of steps and a good policy must generalise to crowds it has never
exactly seen. A shortest-path solver answers ``given this fixed map, what is the quickest
route?''; our question is ``what should I decide now, under uncertainty about what comes
next, to do well on average across many possible days?'' That is a sequential decision
problem and learning is the natural tool for it.

\paragraph{Why does the more advanced method sometimes lose?} Because ``advanced'' is not
``better for this problem.'' Double Q-learning is a genuine improvement on Q-learning
\emph{when there is overestimation to remove} and on the deterministic toy there is none,
so it ties. DQN is a powerful method, but its $\max$ over noisy values is a liability on a
noisy crowd, where PPO's max-free, clipped updates are simply better suited. The competent
move is to match the method to the structure of the problem and our results are a
case study in what happens when you do and do not.

\paragraph{Why average five seeds rather than report the best run?} Because the best run is
cherry-picking and it can hide fragility. The art-lover/max-crowd cell looks like a clean
$+41\%$ on a lucky seed, but across five it is revealed as a coin-flip between a large gain
and a collapse. The mean with its standard deviation is the faithful summary: it states both
the typical outcome and how much we can trust it.

\paragraph{If we hand the agent the legal moves, is it really learning?} Yes, it is learning
the part that matters. Masking removes only the trivial sub-problem of which moves physically
exist; the agent still has to learn \emph{when} to book, \emph{how far} ahead and how to
pace a long visit so every window is met before closing. We deliberately spend its capacity
on those genuine decisions rather than on rediscovering the floor plan.

% =============================================================================
\section{Glossary}
% =============================================================================

\begin{description}[leftmargin=0pt,itemsep=2pt,style=nextline]
\item[State $s$] Everything the agent knows when it must act; a complete-enough snapshot.
\item[Action $a$] A choice available in a state; the legal set can depend on the state.
\item[Reward $r$] The immediate score the world returns after one action.
\item[Return $G$] The (discounted) sum of all rewards from now to the end; what we maximise.
\item[Discount $\gamma$] Weight on future rewards; near $1$ is far-sighted, near $0$ short-sighted.
\item[Policy $\pi$] The strategy: a map from states to actions (or action probabilities).
\item[Value $V(s)$, $Q(s,a)$] Expected return from a state, or from a state-action pair.
\item[MDP / Markov] The formal problem (states, actions, rewards, transitions, discount); ``Markov'' means the state summarises the past well enough that the future depends only on it.
\item[TD error] The gap between a freshly bootstrapped estimate and the old one; the thing every value method shrinks.
\item[Bootstrapping] Updating an estimate toward a target that itself contains an estimate.
\item[On-policy / off-policy] Learning about the policy you are running (PPO) versus about the optimal one while behaving otherwise (Q-learning, DQN).
\item[Overestimation] The upward bias from taking a $\max$ over noisy estimates; cured by Double-Q / Double-DQN.
\item[Function approximation] Using a neural network instead of a table to compute values or probabilities, so similar states share learning.
\item[Gradient descent] Nudging the network's weights in the direction that most reduces a loss.
\item[Advantage $\hat A_t$] How much better an action turned out than the critic expected; the clean signal PPO learns from.
\item[GAE] Generalised advantage estimation; blends short- and long-horizon advantage estimates via $\lambda$.
\item[Clipping] PPO's brake that forbids the policy ratio from moving more than $\pm\epsilon$ per update.
\item[Entropy bonus] A small reward for keeping action probabilities spread out, to sustain exploration.
\item[Action masking] Forcing illegal actions to zero probability before the softmax.
\item[Replay buffer] DQN's memory of past transitions, replayed in random batches.
\item[Target network] A slowly updated copy of the Q-network that supplies stable targets.
\item[Deadly triad] Bootstrapping + function approximation + off-policy learning; together, a recipe for instability.
\item[Common random numbers (CRN)] Scoring every policy on the same fixed crowd scenarios so differences reflect the policy, not luck.
\item[Seed] The random-number starting point of a run; we average over five to get error bars.
\item[Credit assignment] The problem of attributing a delayed outcome to the earlier action that caused it.
\item[Fill curve] The rule tying slot availability to lead time and crowd: busier days need longer leads.
\item[RAMA] Reservation-Anchored Masterpiece Access; the booked-slot system the agent learns to use.
\end{description}

\end{document}
